% ============================================================================
\section{\texttt{int 0x80} 后端实现（D-3）}
\label{sec:int80-backend}
% ============================================================================

在第~\ref{sec:syscall-backend-preview}~节的预告中，我们描述了 \texttt{int~0x80}
后端的设计意图。D-3 将它完整落地：三个新文件协同工作，让 Ring~3 代码能够真正
执行系统调用。

\begin{figure}[H]
  \centering
  \begin{tikzpicture}[
    box/.style={draw, rounded corners=3pt, minimum width=3.5cm,
                minimum height=0.85cm, font=\small, align=center},
    f1/.style={box, fill=blue!8},
    f2/.style={box, fill=red!8},
    f3/.style={box, fill=green!8},
    arr/.style={-{Stealth[length=4pt]}, thick},
    note/.style={font=\scriptsize\itshape, text=gray},
  ]
    \node[f3] (km)   at (0,   0)    {\texttt{kmain.c}\\{\footnotesize 注册后端}};
    \node[f1] (hdr)  at (4.5, 0)    {\texttt{syscall\_int80.h}\\{\footnotesize 声明 get\_ops()}};
    \node[f2] (back) at (0,  -1.8)  {\texttt{syscall\_int80\_backend.c}\\{\footnotesize \texttt{int80\_init()} + ops 表}};
    \node[f2] (asm)  at (4.5,-1.8)  {\texttt{syscall\_int80\_entry.asm}\\{\footnotesize 汇编 stub}};
    \draw[arr] (km)   -- (hdr)  node[above, pos=0.5, note] {\texttt{\#include}};
    \draw[arr] (km)   -- (back) node[left,  pos=0.5, note] {调用 \texttt{get\_ops()}};
    \draw[arr] (back) -- (asm)  node[above, pos=0.5, note] {\texttt{extern} 引用};
    \draw[arr] (hdr)  -- (back) node[right, pos=0.5, note] {\texttt{\#include}};
  \end{tikzpicture}
  \caption{\texttt{int 0x80} 后端的文件组成}
  \label{fig:int80-files}
\end{figure}

% ----------------------------------------------------------------------------
\subsection{IDT 门安装：\texttt{idt\_install\_gate} 与 \texttt{type\_attr}}
\label{sec:int80-idt-gate}
% ----------------------------------------------------------------------------

第一步是在 IDT 的第 \hex{80} 槽安装一个允许 Ring~3 触发的中断门。

\subsubsection{\texttt{type\_attr} 字节：\texttt{0xEE}}

IDT 门描述符中有一个关键的 \texttt{type\_attr} 字节，编码了"谁能触发、进入时是否
关中断"两个属性。\texttt{int 0x80} 使用的值是 \texttt{0xEE}：

\begin{figure}[H]
  \centering
  % fig06-type-attr-bitfield.tex
% type_attr = 0xEE = 1110_1110b 的位域解析图
\begin{tikzpicture}[
  bitcell/.style={
    draw, minimum width=1.0cm, minimum height=0.8cm,
    font=\ttfamily\bfseries\normalsize, align=center,
    anchor=south west,
  },
  bitnum/.style={font=\ttfamily\scriptsize, text=gray!70, align=center},
  lbl/.style={font=\small, align=center},
]

  % ---- 位值（0xEE = 1110 1110）----
  % bit7=1  bit6=1  bit5=1  bit4=0  bit3=1  bit2=1  bit1=1  bit0=0
  % Fields:  P      DPL           S       Type
  \def\bw{1.0}  % bit cell width
  \def\bh{0.8}  % bit cell height

  % bit 7: P=1 (present)
  \node[bitcell, fill=green!20]  at (0*\bw, 0) {1};
  % bit 6-5: DPL=11
  \node[bitcell, fill=blue!15, minimum width=2.0cm] at (1*\bw, 0) {1\ \ 1};
  % bit 4: S=0
  \node[bitcell, fill=yellow!20] at (3*\bw, 0) {0};
  % bit 3-0: Type=1110
  \node[bitcell, fill=orange!15, minimum width=4.0cm] at (4*\bw, 0) {1\ \ 1\ \ 1\ \ 0};

  % ---- 位编号（上方）----
  \foreach \i/\b in {0/7, 1/6, 2/5, 3/4, 4/3, 5/2, 6/1, 7/0} {
    \node[bitnum] at (\i*\bw + 0.5, \bh + 0.25) {\b};
  }
  \node[font=\scriptsize, gray, anchor=east] at (0, \bh+0.25) {bit};

  % ---- 字段标注（下方，带颜色括号）----
  % P
  \draw[decorate, decoration={brace, amplitude=4pt, mirror},
        thick, green!50!black]
    (0*\bw, -0.1) -- (1*\bw, -0.1)
    node[midway, below=8pt, lbl, text=green!50!black]
      {P=1\\{\scriptsize Present}};

  % DPL
  \draw[decorate, decoration={brace, amplitude=4pt, mirror},
        thick, blue!60!black]
    (1*\bw, -0.1) -- (3*\bw, -0.1)
    node[midway, below=8pt, lbl, text=blue!60!black]
      {DPL=3\\{\scriptsize Ring 3 可触发}};

  % S
  \draw[decorate, decoration={brace, amplitude=4pt, mirror},
        thick, yellow!60!black]
    (3*\bw, -0.1) -- (4*\bw, -0.1)
    node[midway, below=8pt, lbl, text=yellow!60!black]
      {S=0\\{\scriptsize 系统门}};

  % Type
  \draw[decorate, decoration={brace, amplitude=4pt, mirror},
        thick, orange!70!black]
    (4*\bw, -0.1) -- (8*\bw, -0.1)
    node[midway, below=8pt, lbl, text=orange!70!black]
      {Type=0xE\\{\scriptsize 32-bit Int Gate}};

  % ---- 总体标注 ----
  \node[font=\bfseries, anchor=south west] at (0, \bh + 0.65)
    {\texttt{type\_attr = 0xEE = 1110\,1110b}};

  % ---- 对比 0x8E ----
  \node[font=\small, anchor=west, text=gray!70] at (8.2, \bh*0.5)
    {cf.\ \texttt{0x8E}:};
  \node[font=\scriptsize, anchor=west, text=gray!60] at (8.2, \bh*0.5 - 0.35)
    {DPL=0（内核专用）};

\end{tikzpicture}

  \caption{\texttt{type\_attr = 0xEE} 的位域解析（DPL=3，32-bit Interrupt Gate）}
  \label{fig:type-attr}
\end{figure}

\begin{description}
  \item[\texttt{bit 7}\enspace P = 1]
    描述符有效（Present）。若 P=0，CPU 触发 \#NP（Segment Not Present）。
  \item[\texttt{bits 6--5}\enspace DPL = 11$_2$ = 3]
    允许 Ring~3 的软件 \texttt{int n} 指令触发此门。若 DPL=0，Ring~3 执行
    \texttt{int 0x80} 时 CPU 检查 $\text{CPL}(3) \leq \text{DPL}(0)$ 不成立，
    直接触发 \#GP。
  \item[\texttt{bit 4}\enspace S = 0]
    系统描述符（Gate 类型），而非代码/数据段。
  \item[\texttt{bits 3--0}\enspace Type = 1110$_2$ = 0xE]
    32-bit Interrupt Gate。进入处理程序时 CPU 自动清除 EFLAGS.IF（关中断），
    防止 syscall 处理期间被硬件中断打断。
\end{description}

对比内核异常/IRQ 使用的 \texttt{0x8E}：两者 Type 字段相同，只有 DPL 不同——
\texttt{0x8E} 的 DPL=0，仅内核态可软件触发。

\begin{warnbox}
\textbf{Interrupt Gate vs Trap Gate}

选择 Interrupt Gate（Type=\texttt{0xE}，进入时清 IF）而非 Trap Gate（Type=\texttt{0xF}，
保持 IF）是出于当前阶段的简化考量。关中断意味着 syscall 处理期间不会被时钟中断
打断，避免了不完整的内核栈被调度器看到的问题。

Linux 的 \texttt{int 0x80} 实际使用 Trap Gate 以支持内核抢占。当本书引入内核
线程调度时，可将 \texttt{0xEE} 改为 \texttt{0xEF} 切换到 Trap Gate。
\end{warnbox}

\subsubsection{\texttt{idt\_install\_gate} 接口}

为了让后端模块在 \texttt{idt\_init()} 完成后仍能动态安装向量，我们在
\texttt{idt.h} 中新增了一个公开接口：

\codefile{src/include/arch/idt.h}
\begin{lstlisting}[style=cstyle]
void idt_install_gate(uint8_t vector, void (*handler)(void),
                      uint16_t selector, uint8_t type_attr);
\end{lstlisting}

实现上只是调用原有的 \texttt{static idt\_set\_gate()}——把它从模块内部提升为可导出
符号。因为 IDTR 指向的是同一片 \texttt{idt[]} 数组，修改表项后 CPU 下次查该向量时
就会用新值，无需重新执行 \texttt{lidt}。

% ----------------------------------------------------------------------------
\subsection{汇编入口 Stub：\texttt{syscall\_int80\_entry}}
\label{sec:int80-asm-stub}
% ----------------------------------------------------------------------------

汇编 stub 是 IDT 门直接跳到的第一段内核代码。它的职责是：
在内核栈上构造 \texttt{syscall\_regs\_t} 结构体 → 切换数据段 → 调用 C 的
dispatch 函数 → 恢复现场 → \texttt{iret} 返回用户态。

\subsubsection{进入时的 CPU 状态}

Ring~3 执行 \texttt{int 0x80} 之后、汇编第一条指令运行之前，CPU 已完成：

\begin{enumerate}
  \item 从 TSS 读取 \texttt{SS0:ESP0}，切换到内核栈（来自第~\ref{ch:tss-ring3}~章配置的 TSS）
  \item 向内核栈依次压入：\texttt{SS\_user} / \texttt{ESP\_user} /
    \texttt{EFLAGS} / \texttt{CS\_user} / \texttt{EIP\_user}（完整跨特权级帧）
  \item 清除 EFLAGS.IF（Interrupt Gate 自动关中断）
  \item 将 CS 设置为 IDT gate 中的 selector（内核代码段 \hex{08}），跳转到 stub
\end{enumerate}

\noindent 注意：DS/ES/FS/GS \textbf{仍然是用户态的值}（\hex{23}），
必须由 stub 手动切换。

\subsubsection{内核栈布局快照}

下图展示了即将调用 \texttt{syscall\_dispatch} 时内核栈的完整状态：

\begin{figure}[H]
  \centering
  % fig05-int80-stack.tex
% 内核栈布局：syscall_int80_entry 调用 syscall_dispatch 前的快照
% 低地址在下（ESP 在最下），高地址在上（iret 帧在最上）
\begin{tikzpicture}[
  cell/.style={
    draw, minimum width=4.8cm, minimum height=0.62cm,
    anchor=south west, font=\ttfamily\small, align=left,
    inner xsep=6pt,
  },
  off/.style={font=\ttfamily\scriptsize, anchor=west, text=gray!70},
  grplbl/.style={font=\small, align=right},
]

  \def\H{0.62}   % cell height (must match minimum height)
  \def\W{4.8}    % cell width

  % ---- 保存的段寄存器（blue）  ESP+0 .. ESP+12 ----
  \node[cell, fill=blue!8]  (gs)   at (0, 0*\H) {gs};
  \node[cell, fill=blue!8]  (fs)   at (0, 1*\H) {fs};
  \node[cell, fill=blue!8]  (es)   at (0, 2*\H) {es};
  \node[cell, fill=blue!8]  (ds)   at (0, 3*\H) {ds};

  % ---- syscall_regs_t（green）  ESP+16 .. ESP+36 ----
  \node[cell, fill=green!8]  (nr)  at (0, 4*\H) {nr \quad\quad\ \ (EAX)};
  \node[cell, fill=green!8]  (a1)  at (0, 5*\H) {arg1 \quad\ \ (EBX)};
  \node[cell, fill=green!8]  (a2)  at (0, 6*\H) {arg2 \quad\ \ (ECX)};
  \node[cell, fill=green!8]  (a3)  at (0, 7*\H) {arg3 \quad\ \ (EDX)};
  \node[cell, fill=green!8]  (a4)  at (0, 8*\H) {arg4 \quad\ \ (ESI)};
  \node[cell, fill=green!8]  (a5)  at (0, 9*\H) {arg5 \quad\ \ (EDI)};

  % ---- CPU iret 帧（orange）  ESP+40 .. ESP+56 ----
  \node[cell, fill=orange!12] (eip)  at (0, 10*\H) {EIP\_user};
  \node[cell, fill=orange!12] (csu)  at (0, 11*\H) {CS\_user \quad\ (\texttt{0x1B})};
  \node[cell, fill=orange!12] (efl)  at (0, 12*\H) {EFLAGS};
  \node[cell, fill=orange!12] (esp3) at (0, 13*\H) {ESP\_user};
  \node[cell, fill=orange!12] (ss3)  at (0, 14*\H) {SS\_user \quad\ (\texttt{0x23})};

  % ---- 地址偏移标注（右侧）----
  \foreach \i/\txt in {
    0/ESP+0,  1/ESP+4,  2/ESP+8,  3/ESP+12,
    4/ESP+16, 5/ESP+20, 6/ESP+24, 7/ESP+28, 8/ESP+32, 9/ESP+36,
    10/ESP+40, 11/ESP+44, 12/ESP+48, 13/ESP+52, 14/ESP+56
  } {
    \node[off] at (\W+0.15, \i*\H + 0.31) {\txt};
  }

  % ---- ESP 箭头（最左侧）----
  \draw[->, thick, red!70!black]
    (-2.1, 0.31) -- (-0.05, 0.31)
    node[left, pos=0, font=\small\bfseries, text=red!70!black] {ESP};

  % ---- &regs = ESP+16 标注 ----
  \draw[->, thick, dashed, blue!70!black]
    (-2.1, 4*\H + 0.31) -- (-0.05, 4*\H + 0.31)
    node[left, pos=0, font=\small, text=blue!70!black]
      {\texttt{\&regs}};
  \node[font=\scriptsize, text=blue!60!black, anchor=east]
    at (-0.1, 4*\H + 0.65)
    {\texttt{lea eax,[esp+16]}};

  % ---- 分组括号（最左侧）----
  % 保存的段
  \draw[decorate, decoration={brace, amplitude=5pt, mirror},
        thick, blue!60!black]
    (-1.95, 0*\H) -- (-1.95, 4*\H)
    node[midway, left=8pt, grplbl, text=blue!60!black]
      {段寄存器\\保存};

  % syscall_regs_t
  \draw[decorate, decoration={brace, amplitude=5pt, mirror},
        thick, green!50!black]
    (-1.95, 4*\H) -- (-1.95, 10*\H)
    node[midway, left=8pt, grplbl, text=green!50!black]
      {\texttt{syscall\_regs\_t}};

  % iret 帧
  \draw[decorate, decoration={brace, amplitude=5pt, mirror},
        thick, orange!70!black]
    (-1.95, 10*\H) -- (-1.95, 15*\H)
    node[midway, left=8pt, grplbl, text=orange!70!black]
      {CPU 自动\\压入\\iret 帧};

  % ---- 地址方向指示 ----
  \node[font=\scriptsize\itshape, gray, anchor=east] at (-2.5, 0.31)  {低地址};
  \node[font=\scriptsize\itshape, gray, anchor=east] at (-2.5, 14*\H + 0.31) {高地址};
  \draw[->, thin, gray!60] (-2.45, 0.9) -- (-2.45, 14*\H - 0.2)
    node[midway, left=2pt, font=\scriptsize\itshape, text=gray!60, rotate=90,
         anchor=south] {地址增长};

\end{tikzpicture}

  \caption{\texttt{syscall\_int80\_entry} 中，\texttt{call syscall\_dispatch}
           前的内核栈布局}
  \label{fig:int80-stack}
\end{figure}

蓝色区域（ESP+0\,\ldots\,+12）保存用户态段寄存器；绿色区域（ESP+16\,\ldots\,+36）
是 \texttt{syscall\_regs\_t} 结构体，\texttt{\&regs = ESP+16}；橙色区域（ESP+40\,\ldots\,+56）
是 CPU 自动压入的 \texttt{iret} 帧。

\subsubsection{完整汇编实现}

\codefile{src/kernel/arch/syscall\_int80\_entry.asm}
\begin{lstlisting}[style=nasmstyle]
extern syscall_dispatch
global syscall_int80_entry

syscall_int80_entry:
    ; 步骤 1：构造 syscall_regs_t（逆序压栈，使结构体从低到高按顺序排列）
    push edi        ; arg5  offset+20
    push esi        ; arg4  offset+16
    push edx        ; arg3  offset+12
    push ecx        ; arg2  offset+8
    push ebx        ; arg1  offset+4
    push eax        ; nr    offset+0  ← ESP 现在指向结构体起始

    ; 步骤 2：保存段寄存器，切换到内核数据段（0x10）
    push ds
    push es
    push fs
    push gs
    mov ax, 0x10    ; GDT_KERNEL_DATA_SEL
    mov ds, ax
    mov es, ax
    mov fs, ax
    mov gs, ax

    ; 步骤 3：调用 syscall_dispatch(&regs)
    ; &regs = ESP+16（段寄存器占了 4×4=16 字节）
    lea eax, [esp + 16]
    push eax                ; 参数：syscall_regs_t*
    call syscall_dispatch
    add esp, 4              ; 清理参数

    ; 步骤 4：恢复段寄存器（DS/ES/FS/GS 回到用户态值）
    pop gs
    pop fs
    pop es
    pop ds

    ; 步骤 5：把 syscall 返回值写入 nr 槽，恢复用户寄存器
    mov [esp], eax  ; 返回值 → nr 槽（pop eax 时取出）
    pop eax         ; EAX = 返回值（用户态通过 EAX 接收结果）
    pop ebx         ; 恢复 arg1
    pop ecx         ; 恢复 arg2
    pop edx         ; 恢复 arg3
    pop esi         ; 恢复 arg4
    pop edi         ; 恢复 arg5

    ; 步骤 6：iret 返回 Ring 3
    iret
\end{lstlisting}

\begin{whybox}
\textbf{为什么用 \texttt{lea eax, [esp+16]} 而不是直接 \texttt{push esp}？}

此时 ESP 指向的是刚压入的 \texttt{gs}（段寄存器保存区的最低地址），
而 \texttt{syscall\_regs\_t} 的起始（\texttt{nr} 字段）在它上方 16 字节处
（4 个段寄存器 × 4 字节）。若直接 \texttt{push esp}，C 函数收到的指针会指向
\texttt{gs}，读出的 \texttt{nr} 值是 DS 选择子，而非 EAX。\texttt{lea} 加上
偏移量才能拿到结构体的真正起始地址。
\end{whybox}

\begin{whybox}
\textbf{为什么要把返回值先写入 \texttt{nr} 槽，再 \texttt{pop eax}？}

\texttt{syscall\_dispatch()} 的返回值存在 EAX 中，但我们需要用一系列 \texttt{pop}
来同步弹出整个 \texttt{syscall\_regs\_t} 并恢复寄存器。第一个 \texttt{pop eax}
会弹出 \texttt{nr} 槽的内容——如果不先覆盖它，弹出的就是系统调用号而非返回值。
\texttt{mov [esp], eax} 把返回值写入 \texttt{nr} 槽，再 \texttt{pop eax} 就能
同时完成"取返回值 + 调整 ESP"两步。
\end{whybox}

\begin{cpustate}
\texttt{iret} 执行后（Ring~0 → Ring~3）：
\begin{itemize}
  \item EIP ← \texttt{EIP\_user}（\texttt{int 0x80} 之后的下一条指令）
  \item CS ← \texttt{CS\_user}（\hex{1B}，DPL=3）→ CPL 变回 3
  \item EFLAGS ← 用户态 EFLAGS（IF 等标志恢复）
  \item ESP ← \texttt{ESP\_user}（用户态栈指针）
  \item SS ← \texttt{SS\_user}（\hex{23}，DPL=3）
  \item EAX = syscall 返回值（用户代码通过 EAX 读取结果）
\end{itemize}
\end{cpustate}

% ----------------------------------------------------------------------------
\subsection{C 后端：\texttt{syscall\_int80\_backend.c}}
\label{sec:int80-c-backend}
% ----------------------------------------------------------------------------

C 后端实现 \texttt{syscall\_ops\_t} 接口，并在 \texttt{init()} 中完成 IDT
门的安装：

\codefile{src/kernel/arch/syscall\_int80\_backend.c}
\begin{lstlisting}[style=cstyle]
/* 0xEE = 1110_1110b: P=1, DPL=11, S=0, Type=1110 (32-bit int gate) */
#define INT_GATE_DPL3  0xEE

extern void syscall_int80_entry(void);   /* 汇编入口（链接时解析） */

static void int80_init(void) {
    idt_install_gate(0x80,
                     syscall_int80_entry,  /* 汇编 stub 地址 */
                     GDT_KERNEL_CODE_SEL,  /* selector = 0x08 */
                     INT_GATE_DPL3);       /* type_attr */

    printk("[syscall/int0x80] IDT[0x80] installed "
           "(DPL=3, selector=0x%02x, entry=%p)\n",
           GDT_KERNEL_CODE_SEL,
           (void *)syscall_int80_entry);
}

static const syscall_ops_t int80_ops = {
    .name = "int0x80",
    .init = int80_init,
};

const syscall_ops_t *syscall_int80_get_ops(void) {
    return &int80_ops;
}
\end{lstlisting}

\begin{whybox}
\textbf{为什么文件名用 \texttt{syscall\_int80\_backend.c}？}

Makefile 用 \texttt{patsubst} 把所有 \texttt{.c} 和 \texttt{.asm} 文件映射到
同名的 \texttt{.o}。若此文件叫 \texttt{syscall\_int80.c}，会与
\texttt{syscall\_int80\_entry.asm} 产生同名目标文件，导致链接重复。
项目中 \texttt{tss.c} + \texttt{tss\_flush.asm} 遵循了相同的命名惯例。
\end{whybox}

% ----------------------------------------------------------------------------
\subsection{\texttt{kmain.c} 集成：后端注册与编译期切换}
\label{sec:int80-kmain}
% ----------------------------------------------------------------------------

\codefile{src/kernel/core/kmain.c}
\begin{lstlisting}[style=cstyle]
/* 头文件包含（编译期选择后端）*/
#if KCONFIG_SYSCALL_BACKEND == 0
#include "syscall_int80.h"
#endif

/* 在 kmain() 中，idt_init() 之后 */
#if KCONFIG_SYSCALL_BACKEND == 0
    syscall_register_backend(syscall_int80_get_ops());
#endif
syscall_init();
\end{lstlisting}

\begin{cpustate}
\texttt{syscall\_init()} 执行后（\texttt{KCONFIG\_SYSCALL\_BACKEND=0}）：
\begin{itemize}
  \item \texttt{IDT[0x80].offset} = \texttt{\&syscall\_int80\_entry}
  \item \texttt{IDT[0x80].selector} = \hex{08}（内核代码段）
  \item \texttt{IDT[0x80].type\_attr} = \hex{EE}（P=1, DPL=3, 32-bit Int Gate）
  \item Ring~3 执行 \texttt{int 0x80}：DPL 检查通过 →
    CPU 切换到 TSS.SS0:ESP0 → 压入 iret 帧 → 跳转到 stub
\end{itemize}
\end{cpustate}

\begin{designbox}
\textbf{\texttt{KCONFIG\_SYSCALL\_BACKEND} 的意义}

\begin{center}
\begin{tabular}{cl}
  \toprule
  值 & 后端 \\
  \midrule
  \texttt{0} & \texttt{int 0x80}（D-3，本节实现） \\
  \texttt{1} & \texttt{sysenter}（D-4，进阶，需配置 MSR） \\
  \bottomrule
\end{tabular}
\end{center}

两个后端共享同一 dispatch 层和 syscall table，切换时只需修改此编译宏并重新
编译——这正是第~\ref{sec:syscall-interface}~节可插拔架构的价值所在。
\end{designbox}

% ----------------------------------------------------------------------------
\subsection{用户态测试代码：PIC 技巧}
\label{sec:int80-user-code}
% ----------------------------------------------------------------------------

\texttt{user\_test\_code.asm} 中新增了 \texttt{user\_syscall\_start/end} 片段，
用于完整验证系统调用路径。

\subsubsection{为什么需要位置无关代码（PIC）}

内核将这段代码当作字节数组复制到用户地址空间（\hex{00400000}），
但汇编时编译器并不知道运行时地址。代码中内嵌了字符串字面量
\texttt{sc\_msg}，若直接用绝对地址引用，复制后地址就失效了。

解决方案是 \textbf{call/pop} 技巧：

\begin{lstlisting}[style=nasmstyle]
    call sc_get_ip   ; 把 sc_get_ip 的地址压栈，然后跳转到 sc_get_ip
sc_get_ip:
    pop  esi         ; ESI = sc_get_ip 的运行时地址
\end{lstlisting}

\texttt{call} 指令执行时，CPU 把紧跟其后的指令地址（即 \texttt{sc\_get\_ip} 标签处）
压入用户栈，然后跳到那里。\texttt{pop esi} 弹出的正是 \texttt{sc\_get\_ip}
的\textbf{运行时}地址。之后用
\texttt{[esi + (sc\_msg - sc\_get\_ip)]} 计算任何其他标签的运行时地址——
括号内的偏移量在编译期已知，无论代码被加载到哪里都保持正确。

\subsubsection{完整测试代码}

\codefile{src/kernel/arch/user\_test\_code.asm}
\begin{lstlisting}[style=nasmstyle]
user_syscall_start:
    ; PIC：获取 sc_get_ip 的运行时地址
    call sc_get_ip
sc_get_ip:
    pop esi                              ; ESI = 运行时 sc_get_ip 地址

    ; write(fd=1, buf=sc_msg, count=18)
    mov eax, 4                           ; SYS_WRITE
    mov ebx, 1                           ; fd = 1 (stdout)
    lea ecx, [esi + (sc_msg - sc_get_ip)]; buf（运行时地址）
    mov edx, sc_msg_len                  ; count
    int 0x80                             ; ← 触发系统调用

    ; exit(status=0)
    mov eax, 1                           ; SYS_EXIT
    xor ebx, ebx                         ; status = 0
    int 0x80                             ; ← 触发系统调用

    jmp $                                ; 安全网（不应执行到此处）

sc_msg:    db "hello from ring3!", 0x0A  ; 18 字节（含换行符）
sc_msg_len equ ($ - sc_msg)
user_syscall_end:
\end{lstlisting}

两次 \texttt{int 0x80} 将分别触发完整的：

\[
\text{Ring~3} \;\xrightarrow{\texttt{int 0x80}}\; \texttt{syscall\_int80\_entry}
  \;\to\; \texttt{syscall\_dispatch} \;\to\; \textit{handler}
\]

第一次调用 \texttt{sys\_write}（nr=4），第二次调用 \texttt{sys\_exit}（nr=1，执行
\texttt{hlt} 令 CPU 停机）。

\begin{thinkbox}
\texttt{call sc\_get\_ip} 会把返回地址压入\textbf{用户栈}，而非内核栈。
\texttt{int 0x80} 触发时 CPU 只保存/恢复 \texttt{ESP\_user}，用户栈上的
这个额外字值不会被内核看到，也不影响后续 \texttt{syscall\_dispatch} 的执行。
为什么用户栈上的内容不需要显式清理？

\emph{提示：\texttt{call sc\_get\_ip} 的 \texttt{call} 指令本身已经跳转到
\texttt{sc\_get\_ip}，它不会"返回"到 \texttt{call} 之后。
用户栈上的返回地址只是作为 PIC 的"临时载体"被 \texttt{pop esi} 消耗掉了。}
\end{thinkbox}
